% Part: first-order-logic
% Chapter: model-theory
% Section: isomorphism

\documentclass[../../../include/open-logic-section]{subfiles}

\begin{document}

\olfileid[pt]{mod}{bas}{iso}
\olsection{Estruturas isomorfas}

As !!{structure}s de primeira ordem podem ser semelhantes de duas
maneiras. Uma delas consiste em tornar verdadeiras as mesmas
!!{sentence}s. Chamamos essas !!{structure}s de \emph{elementarmente
equivalentes}. No entanto, duas estruturas podem ser muito diferentes
e ainda tornar verdadeiras as mesmas !!{sentence}s---por exemplo, uma
pode ser !!{enumerable} e a outra não. Isso ocorre porque há muitas
características de !!a{structure} que não podem ser expressas em
linguagens de primeira ordem, seja porque a linguagem não é rica o
bastante, seja por limitações fundamentais da lógica de primeira ordem,
como o teorema de L\"owenheim--Skolem. Assim, outra maneira, mais
estrita, pela qual duas !!{structure}s podem ser semelhantes é serem
fundamentalmente iguais: diferem apenas nos objetos que as compõem, não
em suas características estruturais. A noção de \emph{isomorfismo}
torna essa ideia precisa.

\begin{defn}
\ollabel{defn:elem-equiv} 
Dadas duas !!{structure}s $\Struct{M}$ e $\Struct M'$ para a mesma
!!{language}~$\Lang{L}$, dizemos que $\Struct{M}$ é
\emph{elementarmente equivalente a} $\Struct M'$, e escrevemos
$\Struct{M} \equiv \Struct M'$, se, e somente se, para toda
!!{sentence}~$!A$ de~$\Lang{L}$, $\Sat{M}{!A}$ se, e somente se,
$\Sat{M'}{!A}$.
\end{defn}

\begin{defn}
\ollabel{defn:isomorphism} Dadas duas !!{structure}s $\Struct{M}$ e
$\Struct M'$ para a mesma !!{language}~$\Lang L$, dizemos que
$\Struct{M}$ é \emph{isomorfa a}~$\Struct M'$, e escrevemos $\Struct{M}
\simeq \Struct M'$, se, e somente se, existe uma função $h \colon
\Domain{M} \to \Domain{M'}$ tal que:
\begin{enumerate}
\item $h$ é !!{injective}: se $h(x) =
  h(y)$, então $x = y$; 
\item $h$ é !!{surjective}: para todo $y \in \Domain{M'}$, existe
  $x \in \Domain{M}$ tal que $h(x) = y$;
\item \ollabel{defn:iso-const}para toda !!{constant} $c$:
  $h(\Assign{c}{M}) = \Assign{c}{M'}$;
\item \ollabel{defn:iso-pred}para todo !!{predicate} de aridade $n$,~$P$:
  \[
  \tuple{a_1, \dots, a_n}\in \Assign{P}{M} \quad\text{se, e somente se,}\quad
  \tuple{h(a_1), \dots, h(a_n)} \in \Assign{P}{M'};
  \]
\item \ollabel{defn:iso-func}para toda !!{function} de aridade $n$, $f$:
  \[
  h(\Assign{f}{M}(a_1, \dots, a_n)) =
  \Assign{f}{M'}(h(a_1), \dots, h(a_n)).
  \]
\end{enumerate}
\end{defn}

\begin{thm}
\ollabel{thm:isom}
Se $\Struct{M} \iso \Struct M'$, então $\Struct{M} \elemequiv
\Struct{M'}$.
\end{thm}

\begin{proof}
Seja $h$ um isomorfismo de $\Struct{M}$ sobre $\Struct M'$. Para toda
atribuição~$s$, $h \circ s$ é a composição de $h$ e $s$, O.s.\@, a
atribuição em $\Struct{M'}$ tal que $(h \circ s)(x) = h(s(x))$. Por
indução em $t$ e $!A$, podem-se provar as afirmações mais fortes:
\begin{enumerate}
  \item[a.] $h(\Value{t}{M}[s]) = \Value{t}{M'}[h\circ s]$.
  \item[b.] $\Sat{M}{!A}[s]$ se, e somente se, $\Sat{M'}{!A}[h \circ s]$.
\end{enumerate}
A primeira é provada por indução na complexidade de~$t$.
\begin{enumerate}
\item Se $t \ident c$, então $\Value{c}{M}[s] = \Assign{c}{M}$ e
  $\Value{c}{M'}[h \circ s] = \Assign{c}{M'}$. Portanto,
  $h(\Value{t}{M}[s]) = h(\Assign{c}{M}) = \Assign{c}{M'}$ (por
  \olref{defn:iso-const} de \olref{defn:isomorphism}) $=
  \Value{t}{M'}[h \circ s]$.
\item Se $t \ident x$, então $\Value{x}{M}[s] = s(x)$ e
  $\Value{x}{M'}[h \circ s] = h(s(x))$. Portanto, $h(\Value{x}{M}[s]) =
  h(s(x)) = \Value{x}{M'}[h \circ s]$.
\item Se $t \ident f(t_1, \dots, t_n)$, então
  \begin{align*}
    \Value{t}{M}[s] & = \Assign{f}{M}(\Value{t_1}{M}[s], \dots, \Value{t_n}{M}[s]) \quad\text{e}\\
  \Value{t}{M'}[h \circ s] & = \Assign{f}{M'}(\Value{t_1}{M'}[h \circ
    s], \dots, \Value{t_n}{M'}[h \circ s]).
  \end{align*}
  A hipótese de indução afirma que, para cada $i$, $h(\Value{t_i}{M}[s])
  = \Value{t_i}{M'}[h\circ s]$. Logo,
  \begin{align}
    h(\Value{t}{M}[s]) 
    & = h(\Assign{f}{M}(\Value{t_1}{M}[s], \dots, \Value{t_n}{M}[s]))\notag\\
    & = \Assign{f}{M'}(h(\Value{t_1}{M}[s]), \dots,
    h(\Value{t_n}{M}[s])) \ollabel{iso-1}\\
    & = \Assign{f}{M'}(\Value{t_1}{M'}[h \circ s], \dots,
    \Value{t_n}{M'}[h \circ s]) \ollabel{iso-2}\\
    & = \Value{t}{M'}[h\circ s] \notag
  \end{align}
  Aqui, \olref{iso-1} decorre de \olref{defn:iso-func} de
  \olref{defn:isomorphism}, e \olref{iso-2}, da hipótese de indução.
\end{enumerate}
A parte (b) fica como exercício.

Se $!A$ é uma sentença, as atribuições~$s$ e $h \circ s$ são
irrelevantes, e temos $\Sat{M}{!A}$ se, e somente se, $\Sat{M'}{!A}$.
\end{proof}

\begin{prob}
Apresente em detalhes a prova de (b) de
\olref[mod][bas][iso]{thm:isom}. Indique em que ponto é usada cada uma
das cinco propriedades que caracterizam os isomorfismos em
\olref[mod][bas][iso]{defn:isomorphism}.
\end{prob}

\begin{defn}
Um \emph{automorfismo} de uma estrutura $\Struct{M}$ é um isomorfismo
de $\Struct{M}$ sobre si mesma.
\end{defn}

\begin{prob}
Mostre que, para toda estrutura $\Struct{M}$, se $X$ é um subconjunto
definível de $\Struct{M}$ e $h$ é um automorfismo de $\Struct{M}$,
então $X = \Setabs{h(x)}{x \in X}$ (O.s.\@, $X$ é fixado por $h$).
\end{prob}


\end{document}
